\clearpage
\VerdierPhysicalPageAnchor{0228}%
\begingroup\small\setlength{\parskip}{0pt}\setlength{\emergencystretch}{1.5em}
\setlength{\abovedisplayskip}{4pt plus 2pt minus 2pt}%
\setlength{\belowdisplayskip}{4pt plus 2pt minus 2pt}%

\VerdierSemanticAnchor{ch4d-4-1-4}
\medskip\noindent\textbf{4.1.4.}\quad
It follows from (4.1.3) that for every object $Y$ of
$\operatorname{comp}(\mathcal A)$,

\VerdierSemanticAnchor{ch4d-4-1-4-spec-k}
\[
\operatorname{Spec}_1(Y)=\bigl(Y(p,q),\delta(p,q,r)\bigr),
\]
defined by (4.1.2.3), (4.1.2.4), and (4.1.2.6), is a spectral object with
values in $K(\mathcal A)$. The spectral object $\operatorname{Spec}_1(Y)$
depends functorially on the object $Y$ of
$\operatorname{comp}(\mathcal A)$. We have therefore obtained a functor

\VerdierSemanticAnchor{ch4d-4-1-4-functor}
\[
Y\longmapsto\operatorname{Spec}_1(Y)
\]
from the category of complexes of $\mathcal A$ to the category of spectral
objects with values in $K(\mathcal A)$. This functor is called the
\emph{first canonical spectral-object functor with values in
$K(\mathcal A)$}. For every object $Y$ of
$\operatorname{comp}(\mathcal A)$, the spectral object
$\operatorname{Spec}_1(Y)$ is called the \emph{first spectral object with
values in $K(\mathcal A)$ associated with the complex $Y$}. Note that when
$Y$ is an object of $\operatorname{comp}^{*}(\mathcal A)$, where $*$ is
$+$, $-$, $b$, or ``empty'', the first spectral object with values in
$K(\mathcal A)$ associated with the complex $Y$ in fact has values in
$K^{*}(\mathcal A)$.

\VerdierSemanticAnchor{ch4d-4-1-5}
\medskip\noindent\textbf{4.1.5.}\quad
Let $\mathcal A$ be an abelian category and let $Y$ be a complex of
$\mathcal A$. The image, under the canonical functor
$K(\mathcal A)\longrightarrow D(\mathcal A)$, of the first spectral object
associated with the complex $Y$ is a spectral object with values in
$D(\mathcal A)$ called the \emph{first spectral object with values in
$D(\mathcal A)$ associated with the complex $Y$}. Most often, we shall
simply say \emph{first spectral object associated with the complex $Y$} and
write (abusively)

\VerdierSemanticAnchor{ch4d-4-1-5-spec-d}
\[
\operatorname{Spec}_1(Y)=\bigl(Y(p,q),\delta(p,q,r)\bigr)
\]
for the first spectral object with values in $D(\mathcal A)$ associated
with the complex $Y$. The composite of the first canonical spectral-object
functor with the canonical functor
$K(\mathcal A)\longrightarrow D(\mathcal A)$ (extended to spectral objects)
is called the \emph{first canonical spectral-object functor with values in
$D(\mathcal A)$}. We emphasize that this is a functor from the category of
complexes of $\mathcal A$ to the category of spectral objects of
$D(\mathcal A)$ and that \emph{this functor does not, in general, factor
through the category of complexes up to homotopy}.

\VerdierSemanticAnchor{ch4d-4-1-6}
\medskip\noindent\textbf{4.1.6.}\quad
$n$-fold complexes of objects of $\mathcal A$ make it possible to construct
simple complexes equipped with filtrations and hence spectral objects with
values in $K(\mathcal A)$. Let $X$ be an $n$-fold complex (Chapter I,
2.1.2), and let $m$ be an integer, $1\leq m\leq n$. For every
$\sigma\in\mathbb Z[n]$ (notation of (Chapter I, 2.1.1)), denote by
$\sigma_m$ its $m$-th coordinate.

\VerdierSemanticAnchor{ch4d-4-1-6-family-setup}
We shall first define a family of complexes and morphisms of complexes. For
every pair $(p,q)$ of $\mathcal{F}\ell_1(\widetilde{\mathbb Z})$, let
${}^{m}X(p,q)$ be the
\endgroup
